\section{Introduction}

\begin{objectivebox}
\begin{itemize}
    \item Model superconductivity as a classical phase transition into mechanical rigidity: phase-locked electron-gears (the AVE analog of Cooper pairs) as meshing topological gyroscopes.
\end{itemize}
\end{objectivebox}

The Meissner Effect describes how superconductors completely expel external magnetic fields. Standard BCS theory models this via the abstract formation of Cooper pairs interacting via phonon exchange.

The AVE framework offers a strictly mechanical interpretation. If fundamental electric charges (electrons) are physical, spinning topological flywheels of the metric network ($0_1$ unknot flux loops carrying circulating helicity), then two phase-locked electron-gears (the AVE analog of a Cooper pair) are simply two such gyroscopes meshing together. In a macroscopic superconducting state or Bose-Einstein Condensate, the thermal vibration (noise) of the lattice drops low enough that the macroscopic population of these individual flywheels structurally interlocks into a single macroscopic, rigid, phase-locked gear-train.

When an external magnetic flux (a boundary torque) attempts to penetrate the superconductor, it physically cannot; turning a single "gear" would require simultaneously driving the combined moment of inertia ($I_{\text{total}}$) of the entire interlocked gyroscope ensemble. Because the total inertia of the phase-locked bulk is large, the boundary gears resist rotation, and this static rejection of boundary torque is the expulsion of the magnetic field. Superconductivity is thus modeled as a classical phase transition into mechanical rigidity. The mechanism, including the resulting London penetration depth, is depicted and derived in Figure~\ref{fig:meissner_gear_train_phase_locked}.

\begin{summarybox}
\begin{itemize}
    \item Superconductivity is modelled as a classical phase transition into mechanical rigidity: phase-locked electron-gears (the AVE analog of Cooper pairs) interlock into a macroscopic gear-train whose large combined moment of inertia rejects boundary torque and expels magnetic flux.
\end{itemize}
\end{summarybox}

\begin{exercisebox}
\begin{enumerate}
    \item Explain how the AVE gear-train model of superconductivity accounts for the Meissner Effect without invoking BCS phonon exchange.
\end{enumerate}
\end{exercisebox}

